\documentclass{article}
\usepackage[utf8]{inputenc}
\usepackage[T1]{fontenc}
\usepackage[margin=2cm]{geometry}
\usepackage{hyperref}
\sloppy
\begin{document}
\section*{Historical source}
Encyclopaedia Britannica scan, source row 1670. Parse 3.5 extracted the text. Super 3.5 VL corrected it from an enlarged crop of each block. It can still contain recognition errors.
Page headers, footers, and page numbers are removed.

\subsection*{Untitled continuation (E0143)}
\textit{Continues an entry from an earlier page.}

Cappado-
cia. lent hands on himself. He is said to have reigned 50
years.
On the death of Archelaus, the kingdom of Cappa-

On the death of Archelaus, the kingdom of Cappadocia was reduced to a Roman province, and governed by those of the equestrian order. It continued subject to the Romans till the invasion of the eastern empire by the Turks, to whom it is now subject, but has no distinguishing modern name. In what was anciently called Cappadocia, however, the Turks have four Beglerbegliks, called Siwas, Trekiand, Marash, and Cogni.

In the time of the Romans, the inhabitants of Cappadocia bore so bad a character, and were reputed so vicious and lewd, that, among the neighbouring nations, a wicked man was emphatically called a Cappadocian. In after ages, however, their lewd disposition was so corrected and refrained by the pure doctrines of Christianity, that no country whatever has produced greater champions of the Christian religion, or given to the church prelates of more unblemished characters.

We have now no system of the Cappadocian laws, and scarce wherewithal to form any particular idea of them. As to their commerce, they carried on a considerable trade in horses, great numbers of which were produced in their country; and we read of them in Scripture as frequenting the fairs of Tyre with this commodity. As Cappadocia abounded with mines of silver, brass, iron, and alum, and afforded great store of slabaster, cryfal, and jasper, it is probable that they might supply the neighbouring countries with these commodities.

The religion of the ancient Cappadocians was much the same with that of the Persians. At Comana there was a rich and stately temple dedicated to Bellona; whose battles the priests and their attendants used to repeat on stated days, cutting and wounding each other as if fired with an enthusiastic fury. No less famous and magnificent were the temples of Apollo Cataenius, and of Jupiter: the last of which had 3000 fixed servants, or religious votaries. The chief priest was next in rank to that of Comana; and, according to Strabo, had an yearly revenue of 15 talents. Diana Perica was worshipped in a city called Capadocia, where women, devoted to the worship of that goddess, were reported to tread barefooted on burning coals, without receiving any hurt. The temples of Diana at Diopolis, and of Annis at Zela, were likewise held in great veneration both by the Cappadocians and Armenians, who flocked to them from all parts. In the latter were tended all oaths in matters of consequence; and the chief among the priests was no way inferior in dignity, power, and wealth, to any in the kingdom; having a royal attendance, and an unlimited authority over all the inferior servants and officers of the temple. The Romans, who willingly adopted all the superstitions and superstitious rites of the nations they conquered, greatly increased the revenues of this and other temples; conferring the priesthood on such as they thought most fit for carrying on their designs. We are told that human sacrifices were offered at Comana; and that this barbarous custom was brought by Orestes and his sister Iphigenia from Taurica Scythica, where men and women were immolated to Diana. But

this custom, if ever it obtained in Cappadocia, was abolished in the times of the Romans.

\subsection*{CAPPANUS (E0144)}
CAPPANUS, a name given by some authors to a worm that adheres to and gnaws the bottoms of ships, to which it is extremely pernicious, especially in the East and West Indies: to prevent this, several ships have lately been sheathed with copper; the first trial of which was made on his majesty's frigate the Alarm. CAPPARIS, in botany: A genus of the monogynia order, belonging to the polyandria class of plants; and in the natural method ranking under the 25th order, Putamine. The calyx is tetraphyllous and coriaceous; there petals are four; the stamina are long; the fruit is a berry, carnous, unilocular, and pedunculated, or furnished with a foot-stalk.

There are sewen species. The fininois, or com- mon copy, is a low limb, generally growing out of the joints of old walls, the lifters of rocks, and a- monglet rabbi, in most of the warm parts of Europe : is in such very little, which lend out many lateral flex- der branches; a matter of that the picked two short crooked liniers, between which and the branches come out the fourth of the leaves, which are single, floors, and Indians around smooth entire list. At letter-mediate joints, between the branches, come out the flower on long footballs, before there endward, the bud with the equipment is gathered for pickling. Thief which are left expand in front of a fingro ross, having five large white points, which are overall and canvary in the middle are placed a great number of long flaminas, surrounding a flyk which rifles above thing, and crowned with an runn germs, which afterwards because our visible filled with kidney-shaped feeds.

Culture. This plant is very differently preferred in Britain; it delights to grow in crevices of rocks, old walls, \&c. and always thrives best in an horizontal posture; so that, when planted either in pots or in the full ground, they seldom thrive, though they may be kept alive for some years. They are propagated by feeds in the warm parts of Europe, but very seldom in Britain.

Ufes. The buds, pickled with vinegar, \&c. are brought to Britain annually from Italy and the Medi-terranean. They are supposed to excite appetite and alist digestion; and to be particularly useful as detergents and aperients in obstructions of the liver and spleen.

\subsection*{CAPRA (E0145)}
CAPRA, or Goat, a genus of quadrupeds belonging to the order of pecora. The horns are hollow, turned upwards, erect, and feabrous. There are eight fore-teeth in the under jaw, and none in the upper; and they have no dog-teeth. This genus consists of 14 species, viz.

I. The mixces, or common goat, with arched carinated horns, and a long beard. It is a native of the cellen mountains.

The goat is an animal of more sagacity than the sheep. Instead of having an antipathy at mankind, they voluntarily mingle with them, and are easily tamed. Even in uninhabited countries, they betray no savage dispositions. In the year 1698, an English vessel having put in to the island of Bonavilla, two negroes came aboard, and offered gratis to the captain as many


\section*{Source record}
\begin{description}
\item[Scan ID] \texttt{britannica:309f765b1ecc83e0a7f7762e44b4a927a7555443:train:1670}
\item[Image file] \texttt{page-1670-ced69857efba.jpg}
\item[Image SHA-256] \texttt{ced69857efbaddd760fb2634b21a1c56dcc9c43fabda3705dd3b69f16e93c707}
\item[Dataset] NationalLibraryOfScotland/encyclopaedia\_britannica\_illustrated at revision \texttt{309f765b1ecc83e0a7f7762e44b4a927a7555443} (CC0-1.0)
\item[Extraction] model \texttt{nvidia/nemotron-3.5-parse}, raw reply \texttt{extract/raw/page-1670-ced69857efba.json}
\item[Raw reply SHA-256] \texttt{8f06e84a904b3de4f1df982ee3a59de6052bfeb8435912e5c773e58a8c2dd331}
\item[Correction] model \texttt{nvidia/Nemotron-3.5-Super-VL}, block crops, raw replies \texttt{correct/raw/crop\_nothink/page-1670-ced69857efba-b*.json}; parse: 1, super\_vl: 13
\item[Running header] G-AP [ 141] C A P
\end{description}
\end{document}
